\section{Prover and Benchmark-Mode Delay}
\label{sec:prover}
To allow for scalability experiments without being heavily constrained by the massive hardware requirements of actual cryptographic proof generation, the prototype relies heavily on a mock proving path for analytical benchmarking, while still supporting a base real prover.

\begin{itemize}
    \item \textbf{Mock Prover Mode:} During standard benchmark execution, the system uses a parameterized time delay to simulate cryptographic overhead. This allows the study to isolate systems-level bottlenecks (sequencer queue growth, database I/O, network latency) from compute-bound cryptographic constraints.
    \item \textbf{Real Prover Backend:} A \texttt{risc0\_prover} implementation exists and is used as the real proof path backend for the executor (\texttt{PROVER\_BACKEND=risc0}). When enabled via \texttt{REQUIRE\_REAL\_PROOFS}, it executes the guest code inside the RISC0 zkVM, generating STARK proofs that can be wrapped in Groth16. It correctly processes the Merkle state updates, exhibiting step-function latency at segment boundaries with a $\sim$70-80s fixed overhead, and produces constant 256-byte proofs.
    \item \textbf{Deprecation of Initial Proposals:} Early progress reports mentioned parallelized proof generation using \texttt{bellman} and \texttt{halo2}. These were deprecated in favor of the current \texttt{risc0} and mock-prover approach to better align with the project's focus on sequencing and batching queue dynamics rather than raw cryptographic circuit optimization.
\end{itemize}
